% ============================================================
% 结果与讨论
% ============================================================
\section{结果与讨论}

\subsection{表面形貌与三维结构}

三张光盘的原始形貌图像见图~\ref{fig:overview}。三组数据均呈现沿慢扫方向延伸的平行条纹，条纹间距离散而规则，与光盘预刻沟槽的螺旋轨道在局部视场内近似为平行直线的预期一致。条纹走向在三张光盘之间不同：cd1 与 cd3 的条纹相对慢扫轴有明显倾斜，cd2 的条纹几乎与慢扫轴平行。原始图像的高度标准差分别为 20.3、19.9 与 31.0~nm，cd3 的起伏最大。

\begin{figure}[H]
\centering
\includegraphics[width=0.95\textwidth]{fig01_overview_raw.pdf}
\caption{三张光盘的原始 AFM 图像（$10\,\mu\mathrm{m}\times10\,\mu\mathrm{m}$，$1024\times1024$ 点）。横向三列为形貌、误差信号（探针起伏）与横向力通道，纵向三行依次为 cd1、cd2、cd3。色标范围取各通道的第 0.5 与 99.5 百分位。}
\label{fig:overview}
\end{figure}

误差信号与横向力通道的图像对比度明显低于形貌通道，但条纹结构清晰可辨，说明三路信号记录的是同一几何结构在不同物理量上的响应。三维形貌（图~\ref{fig:3d}）给出的直观印象是：三张光盘的表面都由“脊—槽”交替构成，脊顶平缓、槽底较窄，与光盘基片上宽岸窄槽的压制结构相符。

\begin{figure}[H]
\centering
\includegraphics[width=0.95\textwidth]{fig05_3d_surface.pdf}
\caption{三张光盘的三维形貌（二阶面整平后，显示时作二倍降采样）。纵轴为高度，单位 nm。}
\label{fig:3d}
\end{figure}

\subsection{处理链的作用与面粗糙度参数}
\label{sec:rough}

以 cd1 为例的处理链见图~\ref{fig:proc}。原始图像中沿 $x$ 方向的整体倾斜与弯曲由两阶处理扣除：逐行一阶拟合扣除了每行自身的倾斜与偏置，二阶面拟合扣除了跨行的慢变面形。处理后沿 $y$ 方向的行均值变化被压到 $10^{-14}$~nm 量级，沿 $x$ 方向的列均值起伏由 73.8~nm 的绝对电平归零。

\begin{figure}[H]
\centering
\includegraphics[width=0.95\textwidth]{fig02_processing_chain.pdf}
\caption{cd1 形貌通道的处理链：(a) 原始图像；(b) 逐行一阶多项式去倾斜；(c) 二阶面整平。三幅图的面粗糙度 $S_q$ 均在 20.26~nm，说明扣除的面形分量不足总方差的 0.1\%。}
\label{fig:proc}
\end{figure}

处理前后的 $S_q$ 几乎不变（20.26~nm），这一结果的物理含义是：三张光盘的表面方差由条纹结构主导，扫描面与被测面之间的整体倾斜与弯曲贡献不足 0.1\%。与此对照，沿 $x$ 方向的列均值在三张光盘上的标准差分别为 8.45、19.60 与 2.20~nm，三者差异并不对应该方向的结构幅度，而是条纹倾斜带来的平均效应——倾斜越大，跨行平均对条纹的抵消越强，见第~\ref{sec:groove} 节。

面粗糙度参数列于表~\ref{tab:rough}。三张光盘的 $S_a$ 在 16.6--27.0~nm、$S_q$ 在 19.9--31.0~nm、峰谷高度 $S_y$ 在 78--116~nm。偏度 $S_{sk}$ 均为负值（$-0.49\sim-0.98$），说明高度分布相对中值偏向低处，即槽的面积占比小于脊；这与表~\ref{tab:groove} 中占宽比大于 0.5 的结果一致。峭度 $S_{ku}$ 为 1.59、2.49 与 1.89，均小于高斯分布的 3，说明高度分布比正态分布更平坦，反映脊顶与槽底各有一段近水平的平台而不是单一峰值的随机起伏。均方根斜率 $S_{dq}$ 与峰密度 $S_{ds}$ 随结构周期缩短而增大：cd3 的周期最短（1090~nm），其 $S_{dq}$ 为 0.223、$S_{ds}$ 为 40.8~$\mu\mathrm{m}^{-2}$，分别是 cd1 的 2.3 倍与 0.46 倍。

\begin{table}[H]
\centering
\caption{三张光盘的面粗糙度参数（二阶面整平后，按 Imager 手册 §47.2 与 ISO 25178 定义计算）}
\label{tab:rough}
\small
\setstretch{1.0}
\begin{tabular}{@{}lccccccc@{}}
\toprule
样品 & $S_a$/nm & $S_q$/nm & $S_y$/nm & $S_z$/nm & $S_{sk}$ & $S_{ku}$ & $S_{dq}$ \\
\midrule
cd1 & 18.50 & 20.26 & 77.9 & 66.5 & $-0.54$ & 1.59 & 0.097 \\
cd2 & 16.61 & 19.94 & 102.4 & 94.5 & $-0.98$ & 2.49 & 0.153 \\
cd3 & 26.95 & 31.02 & 115.9 & 113.5 & $-0.49$ & 1.89 & 0.223 \\
\bottomrule
\end{tabular}
\end{table}

\subsection{沟槽周期与峰谷深度}
\label{sec:groove}

周期由逐行一维傅里叶变换的主峰给出，取 1024 行的中位数，四分位区间作为行间散布的度量。结果列于表~\ref{tab:groove}：cd1 与 cd2 的周期分别为 2170 与 2143~nm，行间四分位区间宽度不超过 15~nm；cd3 的周期为 1090~nm，行间散布仅 3~nm。行间散布小说明三条沟槽结构在视场内的周期高度均匀，测量本身重复性良好。

\begin{table}[H]
\centering
\caption{沟槽几何参数}
\label{tab:groove}
\small
\setstretch{1.0}
\begin{tabular}{@{}lcccccc@{}}
\toprule
样品 & 周期$^\ast$/nm & 行间四分位区间/nm & 条纹倾角/$^\circ$ & 峰谷深度$^\dagger$/nm & 占宽比 & 倾角修正后周期/nm \\
\midrule
cd1 & 2170 & 2167--2175 & $+8.1$ & 55.2 & 0.61 & 2149 \\
cd2 & 2143 & 2137--2151 & $+0.0$ & 60.4 & 0.75 & 2143 \\
cd3 & 1090 & 1089--1092 & $-13.2$ & 93.4 & 0.63 & 1062 \\
\bottomrule
\end{tabular}

\vspace{2pt}
\footnotesize $^\ast$ 逐行一维傅里叶变换主峰的行中位数；$^\dagger$ 逐行 150~nm 平滑剖面的 $P_{99}-P_1$ 中点值。
\end{table}

条纹倾角由式~\eqref{eq:tilt} 给出，cd1 为 $+8.1^\circ$、cd3 为 $-13.2^\circ$、cd2 接近 $0^\circ$。倾角的直接后果体现在平均剖面上：若不做行程对齐，跨 1024 行平均时 cd3 的条纹相位漂移超过半个周期，条纹被完全抹平，平均剖面的峰谷只剩 1.9~nm；对齐后再平均，峰谷恢复到 88.6~nm，与逐行统计的 93.4~nm 吻合。图~\ref{fig:prof} 给出对齐后的平均剖面，cd2 的剖面呈现宽脊与窄槽的矩形轮廓，占宽比 0.75；cd1 与 cd3 的轮廓更接近对称，占宽比 0.61 与 0.63。按占宽比与周期折算，cd1--cd3 的脊宽分别为 1318、1614 与 688~nm，槽宽分别为 852、530 与 403~nm。

\begin{figure}[H]
\centering
\includegraphics[width=0.95\textwidth]{fig03_profiles.pdf}
\caption{三张光盘的平均剖面（去扫描线、面整平、按条纹倾角对齐后沿 $y$ 平均）。细线为对齐后的逐列数据，粗线为 150~nm 平滑结果，箭头标注主周期。}
\label{fig:prof}
\end{figure}

峰谷深度在 55--93~nm 之间，随周期缩短而增大。需要注意的是，这一读数并不等于沟槽的真实深度：探针以有限的曲率半径与锥角接近样品，槽壁对针尖的遮挡使针尖无法到达槽底，测得的槽型被展宽、深度被压缩。手册给出的量级估计为：对台阶，若针尖半锥角为 $10^\circ$，单边展宽约为台阶高度的 18\%；对颗粒，展宽使粒径读数偏大约 20\% 并叠加 $2R$，$R$ 为针尖曲率半径。以此作为量级参照，本文的深度读数应视为下限，真实槽深大于表中数值。

\subsection{频谱特征}

平均剖面的一维功率谱与二维傅里叶谱见图~\ref{fig:spec}。一维谱的主峰波长与表~\ref{tab:groove} 的周期一致，且在 $1/2$、$1/3$ 基频处出现谐波，说明沟槽轮廓接近方波而非正弦——这与三维形貌中脊顶与槽底各有一段平台一致。cd3 的基频峰相对强度最高（其一次谐波与二次谐波清晰可辨），对应其剖面更接近理想周期结构；cd1 的基频峰两侧展宽较大，反映行间周期存在小幅起伏。

二维傅里叶谱的峰位给出条纹法线方向。受视场限制，$k$ 空间的采样间隔为 $0.1\,\mu\mathrm{m}^{-1}$，据此估计的方向角分辨率为数度，与相位斜率法相比低一个量级，故本文的倾角一律取相位斜率法的结果，二维谱只用于确认条纹的一维周期性（能量集中在 $k_x$ 轴上）。

\begin{figure}[H]
\centering
\includegraphics[width=0.95\textwidth]{fig04_spectra.pdf}
\caption{上排：平均剖面的一维功率谱密度（对数坐标，虚线为主周期）；下排：整平后形貌的二维傅里叶谱（对数强度，显示范围 $\pm2\,\mu\mathrm{m}^{-1}$）。}
\label{fig:spec}
\end{figure}

\subsection{横向尺度标定与光盘类型的推断}

三张光盘的周期分别为 2170、2143 与 1090~nm，若直接与格式标准对照，两组都不落在 CD（1600~nm）或 DVD（740~nm）的标称值上。这里存在两种可能：扫描器横向标定偏大，或样品本身周期不同。数据内部提供了区分这两种可能的线索。

第一，cd1 与 cd2 的周期在倾角修正后为 2149 与 2143~nm，相对差仅 0.25\%。两张不同光盘的结构周期一致到这个程度，说明它们属同一规格，且测量条件（同一扫描器、同一视场、同一扫描频率）没有引入可观的行间偏差。

第二，若二者为 CD 类光盘（道距 1600~nm），则扫描横向尺度因子为 $k=2146/1600=1.341$，即仪器标称的 $10\,\mu\mathrm{m}$ 视场实际约为 $13.4\,\mu\mathrm{m}$。以同一因子修正 cd3 得 $1062/1.341=792$~nm，与 DVD 道距 740~nm 相差 $+7.0\%$。

第三，反向假设 cd3 为 DVD 类（$k=1062/740=1.435$），则 cd1 与 cd2 修正为 1496~nm，比 CD 的 1600~nm 低 6.5\%。

两种假设下的残余差异都在 7\% 量级，共同指向的结论是：cd1 与 cd2 为 CD 类光盘、cd3 为 DVD 类光盘，扫描器快速轴的横向标定偏大约 1.34--1.44 倍。7\% 的残余量级与探针在沟槽边缘的卷积、扫描器的非线性以及盘片道距本身的公差相当。

这一结论必须依赖外部验证：本实验没有对扫描器作标准光栅校正，三张光盘的型号也没有记录，因此上文的“标定因子”是以 CD 道距为参照反推的结果，不能作为独立测得的标定值。若要用同一套数据给出无歧义的结论，需要已知周期的标准光栅按手册 §17.1 的流程重做扫描器横向校正，或补记磁盘的实际型号。需要指出的是，横向标定的偏差不影响高度与粗糙度：$z$ 方向由独立的满量程标定给出，与 $x$、$y$ 的伸缩系数无关。

\subsection{误差信号与横向力通道}

三路信号的统计量列于表~\ref{tab:signal}。误差信号的满量程为 1.43~V（cd1、cd2）与 2.00~V（cd3），其标准差折合满量程的 7.6\%、10.3\% 与 13.8\%；横向力的满量程为 0.143~V 与 0.167~V，标准差折合 2.7\%、3.6\% 与 16.9\%。

\begin{table}[H]
\centering
\caption{误差信号与横向力通道的统计量（已扣除逐行偏置）}
\label{tab:signal}
\small
\setstretch{1.0}
\begin{tabular}{@{}llcccc@{}}
\toprule
样品 & 通道 & 满量程 & 标准差 & 峰谷值 & 相对满量程 \\
\midrule
\multirow{2}{*}{cd1} & 误差信号 & 1.429~V & 0.1088~V & 0.943~V & 7.6\% \\
                     & 横向力   & 0.143~V & 0.0039~V & 0.045~V & 2.7\% \\
\multirow{2}{*}{cd2} & 误差信号 & 1.429~V & 0.1478~V & 1.398~V & 10.3\% \\
                     & 横向力   & 0.143~V & 0.0051~V & 0.068~V & 3.6\% \\
\multirow{2}{*}{cd3} & 误差信号 & 2.000~V & 0.2750~V & 1.524~V & 13.8\% \\
                     & 横向力   & 0.167~V & 0.0282~V & 0.154~V & 16.9\% \\
\bottomrule
\end{tabular}
\end{table}

误差信号正比于形貌沿扫描方向的变化率，因而在槽壁处最大、在平台处趋于零。图~\ref{fig:chan} 给出三路信号的单行剖面：形貌通道给出“脊—槽”的绝对高度，误差信号在高度转折处出现尖峰，横向力在槽壁两侧出现符号相反的偏转。后者是摩擦力的典型特征——针尖上坡与下坡时悬臂的扭转变形方向相反，因此横向力信号在每一周期的两侧呈反对称。三张光盘中 cd3 的横向力相对幅度最大（16.9\%），与其结构最深（93.4~nm）和周期最短（1090~nm）一致：在相同的法向载荷下，更陡的槽壁产生更大的扭转。

\begin{figure}[H]
\centering
\includegraphics[width=0.95\textwidth]{fig06_channels.pdf}
\caption{cd1 的三通道对照。上排为原始图像（形貌、误差信号、横向力），下排为第 300、500、700 行的单行剖面。}
\label{fig:chan}
\end{figure}

\subsection{图像质量、伪迹与不确定度}

三组数据均未出现明显的扫描线伪迹。判据是原始图像的行均值序列：其标准差为 0.030、0.033 与 0.056~nm，相邻行之间的跳变小于 0.004~nm，而行内起伏的标准差在 20~nm 量级。若存在反馈失锁或行间机械跳变，行均值会表现出与行内起伏同量级的突变，本数据中没有这种情况。这说明在 1.0~Hz 的扫描频率下，比例与积分增益取 200 时反馈回路对沟槽的高度变化有充分响应。

条纹倾斜是本批数据最需要注意的伪迹。它本身不是测量误差，而是扫描轴与沟槽走向的实际夹角；但若忽略它直接沿慢扫方向平均，剖面幅度会被严重低估（cd3 由 88.6~nm 降至 1.9~nm），进而使“沟槽深度”的估计偏低一半以上。本文用相位斜率法统一估计并校正倾角，三组倾角分别为 $+8.1^\circ$、$0^\circ$ 与 $-13.2^\circ$。

探针卷积是深度读数的系统偏差来源。本文所用探针型号记录为 N/A，无法给出曲率半径与锥角的确切值，因此只能给出量级判断：按手册对台阶的估计（半锥角 $10^\circ$ 时单边展宽约为台阶高度的 18\%），以及槽壁对针尖的几何遮挡，实测深度应系统性小于真实槽深。同类文献中 AFM 测得的 CD 沟槽深度多在一百至二百纳米量级，本文的 55--93~nm 处于该量级的下方，与该判断一致。

成像参数的对照受到数据本身的限制。三组数据在参考点（0.30）、扫描频率（1.0~Hz）、比例与积分增益（200/200）上完全相同，差别只在信号增益（7.0、7.0、5.0）、灵敏度（9.201、9.890、10.118）与满量程高度（105.28、105.28、147.39~nm）。因此本实验能提供的是“信号增益与满量程设置对成像的影響”，而不能单独考察扫描频率与参考点的作用；cd3 的结构最深，但其信号增益反而最低，说明这三组之间的形貌差异来自样品本身，而不是增益设置。若要考察扫描频率与参考点的影响，需要在同一样品上作单变量序列扫描，本实验未做，属本次测量的局限。

% ============================================================
% 结论
% ============================================================
\section{结论}

本文用接触模式原子力显微镜测量了三张光盘的表面形貌，建立了从原始 \texttt{.csm} 文件到粗糙度与沟槽参数的完整处理链，主要结论如下。

（1）三张光盘的表面均由平行沟槽构成，周期分别为 $2170\pm4$、$2143\pm7$ 与 $1090\pm2$~nm，逐行峰谷深度分别为 55.2、60.4 与 93.4~nm，整平后面粗糙度 $S_q$ 为 20.26、19.94 与 31.02~nm。

（2）cd1 与 cd2 的周期在倾角修正后相差 0.25\%，属同一规格；以 CD 道距 1600~nm 为参照反推的扫描横向尺度因子为 1.34，据此 cd3 修正后周期为 792~nm，与 DVD 道距 740~nm 相差 7.0\%。三张光盘判为两张 CD 类与一张 DVD 类，该判断依赖外部型号确认，横向标定值本身不是本实验独立测得。

（3）条纹相对扫描轴存在 $8.1^\circ$ 与 $-13.2^\circ$ 的倾斜，跨行平均前必须作行程对齐；未对齐时平均剖面的峰谷幅度可被低估 45 倍（cd3 由 88.6~nm 降至 1.9~nm）。

（4）三组数据均无扫描线伪迹：行均值标准差不超过 0.06~nm，相邻行跳变小于 0.004~nm，说明在 1.0~Hz、增益 200/200 的条件下反馈回路稳定。

（5）二阶面整平对 $S_q$ 的改变不足 0.1\%，说明表面方差由条纹结构主导；偏度全为负、峭度均小于 3，表明表面为宽脊窄槽的平坦—陡壁交替结构。

（6）深度读数为真值的下限，探针卷积使槽型展宽、深度压缩；本实验探针型号未记录，无法定量修正。误差信号与横向力通道与形貌通道同源，横向力在槽壁两侧呈反对称偏转，cd3 的相对幅度最大，与其结构最深、周期最短一致。

\section{意见与建议}

本次实验在数据采集之外的环节有两处可以改进。第一，扫描器横向标定应作为独立环节先做：用已知周期的标准光栅按手册 §17.1 的流程标定 $X$、$Y$ 快慢扫四个方向，把标定因子作为已知量代入周期测量，而不是用光盘的道距反过来推。这样既能给出绝对尺度，也能暴露扫描器的非线性和各向异性（本机的快扫与慢扫伸缩系数相差 26.8\%，若标定不完整，图像在 $x$、$y$ 方向的尺度会系统不一致）。

第二，样品与探针的信息应当记录。本实验三张光盘的型号与所用探针的曲率半径、锥角均未留下记录，导致周期只能作相对比较、深度无法作卷积修正。建议在采集时就登记盘片类型与探针批次，并在同一样品上补做扫描频率（0.5--3~Hz）与参考点的单变量序列，这对理解“参数对形貌测量的影响”这一实验目标最为直接。

数据处理方面，建议把倾角估计与行列对齐作为标准步骤固定下来。本批数据中若沿用软件默认的“沿慢扫方向取剖面”，两张光盘的沟槽深度会被低估一半以上，而这一错误在剖面上表现为“起伏很小、看不出沟槽”，很容易被误判为扫描失败。

% ============================================================
% 参考文献（待补英文文献后统一排序）
% ============================================================
\section{参考文献}
\renewcommand{\refname}{}
\setstretch{1.25}
\songti\zihao{5}
\begin{thebibliography}{99}

\bibitem{manual_cspm}
本原纳米仪器有限公司. CSPM5500 扫描探针显微镜系统用户手册（V1.11.10 中文版）[Z]. 北京: Being Nano-Instruments, 2013.

\bibitem{mironov2004}
Mironov V L. Fundamentals of Scanning Probe Microscopy[M]. Moscow: The Russian Academy of Sciences, 2004.

\bibitem{binnig1982}
Binnig G, Rohrer H, Gerber C, et al. Surface studies by scanning tunneling microscopy[J]. Physical Review Letters, 1982, 49(1): 57--61.

\bibitem{binnig1986}
Binnig G, Quate C F, Gerber C. Atomic force microscope[J]. Physical Review Letters, 1986, 56(9): 930--933.

\bibitem{sun2005}
孙大许, 刘万里, 马强, 等. 基于 AFM 的光盘形貌研究[J]. 光学仪器, 2005(5).

\bibitem{xu2008}
徐琳, 李伯全, 丁建宁, 等. 基于原子力显微镜的光盘表面微结构的检测[J]. 光学仪器, 2008(4).

\bibitem{jing2003}
景蔚萱, 蒋庄德. 原子力显微镜（AFM）在光盘检测及其质量控制中的应用[J]. 光学精密工程, 2003(4).

\end{thebibliography}

% ============================================================
% 附录
% ============================================================
\newpage
\appendix
\section{数据处理流程与代码清单}

本文的全部数据处理以 Python 实现，不依赖 Windows 后处理软件。各脚本的功能与对应关系列于表~\ref{tab:code}，处理流程为：

\begin{enumerate}
  \item 解析 \texttt{.csm}：读取 BMP 头与像素区，按式~\eqref{eq:decode} 还原 16 位量并换算为 nm 或 V，解码末尾 GBK 参数块；
  \item 逐行一阶多项式拟合减除，消除行倾斜与行偏置；
  \item 全局二阶二维多项式拟合减除，完成面整平；
  \item 逐行一维傅里叶变换取主峰周期，取 1024 行的中位数与四分位区间；
  \item 逐行离散傅里叶变换的相位斜率法估计条纹倾角，按行插值平移后沿 $y$ 平均得平均剖面；
  \item 在平均剖面上取 $P_{99}-P_1$ 与占宽比，在整平后图像上按 ISO 25178 计算面粗糙度；
  \item 计算一维功率谱与二维傅里叶谱，作三维形貌与各通道剖面图。
\end{enumerate}

\begin{table}[H]
\centering
\caption{数据处理脚本清单}
\label{tab:code}
\small
\setstretch{1.0}
\begin{tabularx}{\textwidth}{@{}lX@{}}
\toprule
脚本 & 功能 \\
\midrule
\texttt{csmio.py} & \texttt{.csm} 格式解析：BMP 头、像素反编码、参数块解码 \\
\texttt{afmproc.py} & 处理算子：逐行去倾斜、面整平、粗糙度、周期与倾角、功率谱 \\
\texttt{analyze.py} & 主流程：批量处理九个数据文件，输出参数表与数组 \\
\texttt{plot.py} & 出图：原始图像、处理链、剖面、频谱、三维形貌、通道对照 \\
\bottomrule
\end{tabularx}
\end{table}

所有长度单位为 nm，高度由文件头的满量程高度换算，误差与横向力通道由满量程电压换算。原始数据文件、脚本与结果表一并在实验目录中归档，可复现本文全部数值。

\end{document}
